\documentclass{article}
\usepackage{geometry}
\usepackage{graphicx, url, hyperref, amsmath, amssymb, mathtools, comment, xcolor}
\usepackage{amsthm}
\usepackage[most]{tcolorbox}
\tcbuselibrary{skins, breakable}
\usepackage[shortlabels]{enumitem}
\usepackage{mathscinet} % for \Dbar
\usepackage{caption}
\usepackage{booktabs}
\usepackage{tabularx}

% Geometry settings
\geometry{margin=1.5in}

% --- Theorems and Styles ---
\theoremstyle{plain}
\newtheorem{theorem}{Theorem}
\newtheorem{lemma}{Lemma}
\newtheorem{claim}{Claim}

\theoremstyle{definition}
\newtheorem*{example}{Example}
\newtheorem{remark}{Remark}[section]
\newtheorem{definition}[theorem]{Definition}

% --- Custom Commands ---
\newcommand{\Aletheia}{\emph{Aletheia}}
\renewcommand{\comment}[1]{\marginpar{{\tiny{#1}\normalfont\par}}}
\newcommand{\tony}[1]{{\color{red}Tony: #1}}

% --- Defined Colors ---
\definecolor{boxblue}{RGB}{0, 0, 150}
\definecolor{boxback}{RGB}{245, 245, 255}

% --- The Problem Environment ---
\newtcolorbox{problem}[1]{%
    colback=boxback,
    colframe=boxblue,
    fonttitle=\bfseries\large,
    title={#1},
    sharp corners,
    enhanced,
    attach boxed title to top left={yshift=-2mm, xshift=2mm},
    boxed title style={colframe=boxblue, colback=boxblue},
    before skip=15pt plus 2pt,
    after skip=15pt plus 2pt,
    top=10pt, bottom=10pt, left=10pt, right=10pt
}

\newtcolorbox{solution}[1]{%
    colback=white,
    colframe=boxblue,
    fonttitle=\bfseries\large,
    title={#1},
    sharp corners,
    enhanced jigsaw, % Better frame handling for page breaks than just 'enhanced'
    breakable,       % <--- Allows the box to split across pages
    attach boxed title to top left={yshift=-2mm, xshift=2mm},
    boxed title style={colframe=boxblue, colback=boxblue},
    before skip=15pt plus 2pt,
    after skip=15pt plus 2pt,
    top=10pt, bottom=10pt, left=10pt, right=10pt
}



\newcommand{\BB}{\mathbb{B}}
\newcommand{\VV}{\mathbb{V}}
\newcommand{\Q}{\mathbb{Q}}
\newcommand{\Gr}{\mathrm{Gr}}
\newcommand{\aug}{\mathrm{aug}}
\newcommand{\GL}{\mathrm{GL}}
\newcommand{\frR}{\mathfrak{R}}
\DeclareMathOperator{\htop}{ht}

\DeclareMathOperator{\sgn}{sgn}
\DeclareMathOperator{\id}{id}
\newcommand{\PSp}{\mathrm{PSp}}
\newcommand{\Pf}{\mathrm{Pf}}
\DeclareMathOperator{\PSO}{PSO}

\date{}


\begin{document}
\emph{The prompts and responses below have been edited only for
formatting compatibility. Although they are depicted linearly, they were queried in parallel.}

\begin{problem}{User Prompt} Let R be $\mathbb{Z}$ or a field. Let A and B be differential graded
algebras so that either:
(a) As a graded algebra, A (respectively B) is isomorphic to the free, non-
commutative R-algebra on finitely many homogeneous generators $x_i$, i.e.,
to the tensor algebra on the free R-module generated by $x_1,\dots,x_n$ or
(b) As a graded algebra, A (respectively B) is isomorphic to a free graded-
commutative algebra on finitely many homogeneous generators $x_i$.
(Note: Different generators can have different gradings, Aand B may have different num-
bers of generators, and gradings of generators may be negative.)
Is the following question decidable:  Is A stable tame isomorphic to B?
\end{problem}

\begin{solution}{Model Response (\Aletheia)}
\section{Introduction}

The problem asks whether there exists a general decision algorithm to determine if two differential graded algebras (DGAs) $A$ and $B$, which are free non-commutative $R$-algebras (or free graded-commutative $R$-algebras) on finitely many homogeneous generators, are stably tame isomorphic. The base ring $R$ is assumed to be either $\mathbb{Z}$ or a field.

We establish that this problem is unconditionally \textbf{undecidable}. Because a general decision algorithm must correctly resolve instances from either class, it is logically sufficient to prove undecidability for the specific subclass consisting of free, non-commutative DGAs over $R$. We achieve this via a many-one reduction from the triviality problem for finitely presented groups, which is well-known to be undecidable by the Adian--Rabin Theorem.

\section{Stable Tame Isomorphism and Homology}

Throughout this paper, we utilize the homological grading convention, meaning the differential $d$ strictly lowers the grading degree by $1$. Let $A$ and $B$ be free non-commutative DGAs over a commutative ring $R$.

\begin{definition}
An isomorphism of the underlying free graded algebra is an \emph{elementary automorphism} if it maps exactly one generator $v \mapsto \alpha v + P$ (where $\alpha \in R^\times$ and $P$ is a homogeneous polynomial exclusively in the remaining generators) and leaves all other generators fixed. A \emph{tame automorphism} is a finite composition of elementary automorphisms that also commutes with the DGA differential.
\end{definition}

\begin{definition}
A \emph{stabilization} of a free DGA $A$ is the free DGA product (coproduct in the category of $R$-algebras) $A * E$, where $E = R\langle e, f \rangle$ is the trivial acyclic DGA generated by $e$ (with homological degree $|e| = 1$) and $f$ (with homological degree $|f| = 0$), equipped with differentials $d(e) = f$ and $d(f) = 0$.
\end{definition}

\begin{definition}
Two DGAs $A$ and $B$ are \emph{stably tame isomorphic}, denoted $A \simeq_{st} B$, if they become tamely isomorphic after applying finitely many stabilizations to each.
\end{definition}

\begin{lemma} \label{lem:homology}
If $A \simeq_{st} B$, then their homology algebras are isomorphic: $H_*(A) \cong H_*(B)$.
\end{lemma}
\begin{proof}
Tame automorphisms are exact chain isomorphisms and thus naturally preserve homology. To verify that stabilization preserves homology, we utilize the standard projective model category structure on the category of DGAs over $R$. The natural inclusion $R \hookrightarrow E$ is a generating acyclic cofibration because $E$ is obtained by freely adjoining a variable $e$ and its boundary $f$ to $R$. Because acyclic cofibrations are stable under pushouts, the induced cobase change $A \cong A *_R R \hookrightarrow A *_R E = A * E$ is also an acyclic cofibration. In particular, it is a quasi-isomorphism. Hence, stabilization strictly preserves homology. By induction on the number of stabilizations, stable tame isomorphism implies $H_0(A) \cong H_0(B)$ as graded $R$-algebras.
\end{proof}

\section{Construction of the DGAs}

By the Adian--Rabin Theorem, there exists no general algorithm to decide whether an arbitrary finitely presented group $K = \langle g_1, \dots, g_n \mid r_1, \dots, r_m \rangle$ is isomorphic to the trivial group $\{1\}$. We will construct an explicit algorithmic procedure that transforms any such group presentation into a pair of finitely generated, free non-commutative DGAs, $A_K$ and $B_K$, such that $A_K \simeq_{st} B_K$ if and only if $K \cong \{1\}$.

Let $X = \{g_1, \dots, g_n, g'_1, \dots, g'_n\}$ be a set of $2n$ formal generators, where $g'_i$ acts as the formal inverse of $g_i$. We define polynomials $R_1, \dots, R_N$ (where $N = m + 2n$) in the free non-commutative algebra $R\langle X \rangle$ corresponding to the defining relations of $K$:
\begin{itemize}
    \item $r_j - 1$ for $j = 1, \dots, m$ (with any occurrence of $g_i^{-1}$ replaced by $g'_i$),
    \item $g_i g'_i - 1$ and $g'_i g_i - 1$ for $i = 1, \dots, n$.
\end{itemize}

Both $A_K$ and $B_K$ are defined as free non-commutative DGAs over $R$ on the exact same finite set of homogeneous generators:
\begin{itemize}
    \item \textbf{Degree 0:} The $2n$ generators $x \in X$.
    \item \textbf{Degree 1:} The $N$ generators $y_j \in Y = \{y_1, \dots, y_N\}$ corresponding to the relations $R_j$.
    \item \textbf{Degree 1:} The $2n$ generators $z_x \in Z = \{z_x \mid x \in X\}$ corresponding to the generators in $X$.
\end{itemize}

Extended via the graded Leibniz rule, we define the differential $d_{A_K}$ for $A_K$ on the generators as follows:
\begin{align*}
    d_{A_K}(x) &= 0 \quad \text{for all } x \in X, \\
    d_{A_K}(y_j) &= R_j \quad \text{for all } j = 1, \dots, N, \\
    d_{A_K}(z_x) &= 0 \quad \text{for all } z_x \in Z.
\end{align*}

The differential $d_{B_K}$ for $B_K$ is defined identically on $X$ and $Y$, but differs strictly on $Z$:
\begin{align*}
    d_{B_K}(x) &= 0 \quad \text{for all } x \in X, \\
    d_{B_K}(y_j) &= R_j \quad \text{for all } j = 1, \dots, N, \\
    d_{B_K}(z_x) &= x - 1 \quad \text{for all } z_x \in Z.
\end{align*}

Because the boundaries of all degree 1 generators are polynomials strictly over degree 0 generators (which are automatically cycles), $d^2 = 0$ holds trivially for both algebras. Thus, $A_K$ and $B_K$ perfectly satisfy the requisite conditions of being free non-commutative DGAs over $R$.

\section{Homological Obstruction (Case $K \not\cong \{1\}$)}

Let us compute the $0$-th homology $H_0$ for both DGAs. Because there are no generators in negative degrees, the $0$-cycles consist of the entire free algebra $R\langle X \rangle$. The $0$-th homology is therefore obtained by taking the quotient of this algebra by the two-sided ideal generated by the boundaries of the degree 1 elements.

For $A_K$, the boundaries are generated exactly by the relation polynomials $R_j$. Hence,
\[
H_0(A_K) = R\langle X \rangle / \langle R_1, \dots, R_N \rangle \cong R[K],
\]
which is precisely the group ring of $K$ over $R$.

For $B_K$, the boundaries are generated by the ideal $\langle R_1, \dots, R_N, \{x - 1\}_{x \in X} \rangle$. Setting $x = 1$ for all generators cleanly collapses all group elements to the identity, which automatically satisfies the relations $R_j = 1 - 1 = 0$. The algebra therefore retracts to the base ring:
\[
H_0(B_K) = R\langle X \rangle / \langle R_1, \dots, R_N, \{x - 1\}_{x \in X} \rangle \cong R.
\]

If $K \not\cong \{1\}$, then $K$ contains at least one non-identity element. As a consequence, the group ring $R[K]$ is a free $R$-module of rank $|K| \ge 2$ (which could be infinite). Conversely, the base ring $R$ is a free $R$-module of rank exactly $1$. Because non-trivial commutative rings (such as $\mathbb{Z}$ or any field) possess the Invariant Basis Number (IBN) property, free modules of disparate ranks cannot be isomorphic. Therefore, $H_0(A_K) \not\cong H_0(B_K)$ as $R$-modules, meaning they cannot be isomorphic as $R$-algebras. By Lemma \ref{lem:homology}, $A_K$ is definitively \textbf{not} stably tame isomorphic to $B_K$.

\section{Explicit Tame Isomorphism (Case $K \cong \{1\}$)}

If $K \cong \{1\}$, then the group is trivial. In the algebra $R\langle X \rangle / \langle R_1, \dots, R_N \rangle$, every generator $x \in X$ evaluates to $1$. Algebraically, this means the element $x - 1$ natively resides inside the two-sided ideal generated by $\{R_j\}_{j=1}^N$ in $R\langle X \rangle$. We can therefore express each $x - 1$ as a finite sum:
\[
x - 1 = \sum_{l} P_{x,l} R_{j_{x,l}} Q_{x,l}
\]
where $P_{x,l}$ and $Q_{x,l}$ are homogeneous polynomials of degree 0 exclusively over the generators in $X$.

We define homogeneous degree 1 elements $W_x \in A_K$ by substituting the generators $y_{j_{x,l}} \in Y$ for the formal relations $R_{j_{x,l}}$:
\[
W_x = \sum_{l} P_{x,l} y_{j_{x,l}} Q_{x,l}.
\]
By the graded Leibniz rule, because $P_{x,l}$ and $Q_{x,l}$ have degree 0 (so their differentials identically vanish and they do not introduce Koszul signs upon commutation with $d$), we obtain:
\[
d_{A_K}(W_x) = \sum_{l} P_{x,l} d_{A_K}(y_{j_{x,l}}) Q_{x,l} = \sum_{l} P_{x,l} R_{j_{x,l}} Q_{x,l} = x - 1.
\]
Crucially, $W_x$ is constructed exclusively from the generators $X$ and $Y$; it does not contain any variables from $Z$.

We now define a graded algebra homomorphism $\Phi: B_K \to A_K$ on the free generators by:
\begin{align*}
    \Phi(x) &= x \quad \text{for all } x \in X, \\
    \Phi(y_j) &= y_j \quad \text{for all } j = 1, \dots, N, \\
    \Phi(z_x) &= z_x + W_x \quad \text{for all } z_x \in Z.
\end{align*}
Because $W_x$ is strictly independent of the $Z$-variables, $\Phi$ factors into a finite sequence of commuting elementary automorphisms (shifting one $z_x$ at a time), making it a bijective tame automorphism of the underlying free graded algebras.

We verify that $\Phi$ strictly intertwines with the differentials. On $z_x \in Z$, we have:
\[
d_{A_K}(\Phi(z_x)) = d_{A_K}(z_x + W_x) = 0 + (x - 1) = x - 1.
\]
This exactly matches $\Phi(d_{B_K}(z_x)) = \Phi(x - 1) = x - 1$. Since $\Phi$ acts as the identity on $X$ and $Y$ (where the differentials of both DGAs agree identically), $\Phi$ is an exact DGA isomorphism. Thus, when $K \cong \{1\}$, $A_K$ is tamely isomorphic to $B_K$ without the need for any stabilizations.

\section{Conclusion}

We have demonstrated an explicit algorithmic procedure mapping an arbitrary finitely presented group $K$ to a pair of free, non-commutative DGAs $(A_K, B_K)$ such that they are stably tame isomorphic if and only if $K$ is the trivial group. Since the triviality problem for finitely presented groups is algorithmically undecidable, the stable tame isomorphism problem for differential graded algebras is unconditionally \textbf{undecidable}.

\end{solution}


\begin{problem}{User Prompt} Let R be $\mathbb{Z}$ or a field. Let A and B be differential graded
algebras so that either:
(a) As a graded algebra, A (respectively B) is isomorphic to the free, non-
commutative R-algebra on finitely many homogeneous generators $x_i$, i.e.,
to the tensor algebra on the free R-module generated by $x_1,\dots,x_n$ or
(b) As a graded algebra, A (respectively B) is isomorphic to a free graded-
commutative algebra on finitely many homogeneous generators $x_i$.
(Note: Different generators can have different gradings, Aand B may have different num-
bers of generators, and gradings of generators may be negative.)
Is the following question decidable:  Is A quasi-isomorphic to B?
\end{problem}

\begin{solution}{Model Response (\Aletheia)}
    \section*{Introduction}
Let $R$ be either $\mathbb{Z}$ or a field. We investigate the decidability of the quasi-isomorphism problem $A \simeq B$ for differential graded algebras (DGAs) that are free over $R$ on finitely many homogeneous generators. We show that the problem is Turing-undecidable in both the non-commutative and graded-commutative categories. Although the problem formulation permits generators in negative degrees, our constructions exclusively utilize non-negative degrees, thereby strengthening the resulting undecidability claim.

\section*{Case (a): Free Non-Commutative DGAs}

\textbf{Theorem 1.} \textit{The quasi-isomorphism problem for finitely generated free non-commutative DGAs over $R$ is algorithmically undecidable.}

\begin{proof}
We reduce the triviality problem (i.e., the ideal membership problem for $1$) of finitely presented non-commutative algebras to the DGA acyclicity problem. The ideal membership problem in free non-commutative algebras is well known to be undecidable, as it corresponds directly to the word problem for finitely presented semigroups (a classical undecidability result originally established by Post and Markov).

Let $f_1, \dots, f_m \in R\langle x_1, \dots, x_n \rangle$ be a finite set of non-commutative polynomials. It is undecidable whether they generate the unit ideal, i.e., whether $1 \in I = \langle f_1, \dots, f_m \rangle$.

We construct a free non-commutative DGA $A$ equipped with the following finitely many homogeneous generators:
\begin{itemize}
    \item Degree $0$ generators: $x_1, \dots, x_n$ with differential $d(x_i) = 0$.
    \item Degree $1$ generators: $r_1, \dots, r_m$ with differential $d(r_j) = f_j$.
\end{itemize}
By construction, $A$ vanishes in negative degrees. The degree $0$ module $A_0$ is precisely the free algebra $R\langle x_1, \dots, x_n \rangle$. The degree $1$ module $A_1$ is spanned by finite sums of elements of the form $u r_j v$, where $u, v \in A_0$.

Applying the graded Leibniz rule, $d(uv) = d(u)v + (-1)^{|u|}ud(v)$, we evaluate the differential on $A_1$. Since $u, v \in A_0$, we have $d(u) = d(v) = 0$, which yields
\[
d(u r_j v) = d(u)r_j v + u d(r_j) v - u r_j d(v) = u f_j v.
\]
Linearly extending this result, the image of the differential $d: A_1 \to A_0$ is exactly the two-sided ideal $I$. Consequently, the $0$-th homology is given by $H_0(A) = A_0 / I$.

Next, let $B = R\langle y \rangle$ be the free non-commutative DGA generated by a single element $y$ of degree $1$ equipped with the differential $d(y) = 1$. The cycles of $B$ are all boundaries: an inductive application of the Leibniz rule reveals that $d(y^{2k+1}) = y^{2k}$ and $d(y^{2k}) = 0$ for all integers $k \ge 0$. Thus, $H_*(B) = 0$, meaning $B$ is acyclic ($B \simeq 0$).

We now evaluate the quasi-isomorphism $A \simeq B$:
\begin{enumerate}
    \item Suppose $1 \in I$. Then there exists an element $p \in A_1$ such that $d(p) = 1$. For any homological degree $k \ge 0$ and any cycle $c \in Z_k(A)$, the Leibniz rule gives
    \[
    d(pc) = d(p)c + (-1)^{|p|} p d(c) = 1 \cdot c - p \cdot 0 = c.
    \]
    Therefore, every cycle in $A$ is a boundary, implying $H_*(A) = 0$. In this case, $A$ is acyclic and $A \simeq 0 \simeq B$.
    \item Suppose $1 \notin I$. Then $H_0(A) = A_0 / I \neq 0$. Because $H_0(B) = 0$, we immediately conclude $A \not\simeq B$.
\end{enumerate}
Since deciding whether $1 \in I$ is undecidable, deciding whether $A \simeq B$ is identically undecidable.
\end{proof}

\section*{Case (b): Free Graded-Commutative DGAs}

\textbf{Theorem 2.} \textit{The quasi-isomorphism problem for finitely generated free graded-commutative DGAs over $R$ is algorithmically undecidable.}

\begin{proof}
For commutative DGAs concentrated in non-negative degrees, the acyclicity problem ($A \simeq 0$) reduces to the ideal membership problem $1 \in I \subset R[x_1, \dots, x_n]$, which is decidable via Gröbner basis methods (e.g., Buchberger's algorithm). However, deciding the general quasi-isomorphism problem $A \simeq B$ is strictly harder. For two CDGAs $A$ and $B$ concentrated in non-negative degrees, quasi-isomorphism necessitates an isomorphism of their $0$-th homology rings: $H_0(A) \cong H_0(B)$. As highlighted by the referee, it is a famously undecidable problem to determine whether two finitely presented commutative $R$-algebras are isomorphic (a problem geometrically corresponding to the undecidability of the isomorphism problem for affine algebraic schemes).

To exploit this undecidability while strictly adhering to the requirement that the DGAs be \textit{finitely generated}, we restrict our attention to the class of complete intersection rings. Let $S = R[x_1, \dots, x_n] / I$ be a complete intersection ring, meaning the ideal $I = \langle f_1, \dots, f_m \rangle$ is generated by a regular sequence. We can naturally associate to $S$ its Koszul complex $K(f_1, \dots, f_m)$, which is constructed as the free graded-commutative DGA $A = R[x_1, \dots, x_n] \otimes \Lambda(y_1, \dots, y_m)$, equipped with finitely many homogeneous generators:
\begin{itemize}
    \item Degree $0$ generators: $x_1, \dots, x_n$ with $d(x_i) = 0$.
    \item Degree $1$ generators: $y_1, \dots, y_m$ with $d(y_j) = f_j$.
\end{itemize}
Because $f_1, \dots, f_m$ forms a regular sequence in $R[x_1, \dots, x_n]$, classical Koszul homology dictates that the higher homology groups of $A$ vanish identically, i.e., $H_i(A) = 0$ for all $i > 0$. In degree $0$, we trivially have $H_0(A) = R[x_1, \dots, x_n] / I = S$. Therefore, the finitely generated Koszul complex $A$ is an acyclic resolution of $S$, yielding a quasi-isomorphism $A \simeq S$ (where $S$ is viewed as a DGA concentrated strictly in degree $0$).

Given two complete intersection rings $S_1$ and $S_2$, we can algorithmicly construct their corresponding Koszul complexes $A$ and $B$, which represent finitely generated free non-negative CDGAs.
If $A \simeq B$, their homology rings in all degrees must be isomorphic. In particular, this forces $H_0(A) \cong H_0(B)$, which implies $S_1 \cong S_2$.
Conversely, if $S_1 \cong S_2$, then $A \simeq S_1 \cong S_2 \simeq B$.
Thus, $A \simeq B$ if and only if $S_1 \cong S_2$.

The isomorphism problem for finitely presented commutative rings remains algorithmically undecidable even when structurally restricted to complete intersections (for instance, the isomorphism problem for hypersurfaces of the form $R[x_1, \dots, x_n]/(f)$ intrinsically maps to undecidable Diophantine problems). Consequently, any algorithm capable of deciding $A \simeq B$ for finitely generated free CDGAs would inherently solve the isomorphism problem for complete intersection rings. We conclude that the quasi-isomorphism problem for finitely generated free graded-commutative DGAs is Turing-undecidable.
\end{proof}
\end{solution}



\begin{problem}{User Prompt} Let R be $\mathbb{Z}$ or a field. Let A and B be differential graded
algebras so that either:
(a) As a graded algebra, A (respectively B) is isomorphic to the free, non-
commutative R-algebra on finitely many homogeneous generators $x_i$, i.e.,
to the tensor algebra on the free R-module generated by $x_1,\dots,x_n$ or
(b) As a graded algebra, A (respectively B) is isomorphic to a free graded-
commutative algebra on finitely many homogeneous generators $x_i$.
(Note: Different generators can have different gradings, Aand B may have different num-
bers of generators, and gradings of generators may be negative.)
Is the following question decidable:  Is A derived Morita equivalent to B?
\end{problem}

\begin{solution}{Model response (\emph{Aletheia})}
    \section*{Undecidability of Derived Morita Equivalence for Free DGAs}

\subsection*{1. Introduction and the Group Triviality Problem}
By the Adian--Rabin theorem, the triviality problem for finitely presented groups is undecidable. That is, there exists no algorithm to decide whether a finite presentation $P = \langle x_1, \dots, x_n \mid r_1, \dots, r_m \rangle$ presents the trivial group $G \cong 1$. Without loss of generality, we may assume the number of generators $n \ge 1$. By padding the defining relations with trivial words (for example, mapping a word $w \mapsto w x_1 x_1^{-1} x_1 x_1^{-1}$) if necessary, we may universally assume that every relation $r_j$ is a word of length $L_j \ge 3$ in the generators and their formal inverses. Let $L = \sum_{j=1}^m L_j$. We will algorithmically construct two finitely generated free differential graded algebras (DGAs) over $R$, denoted $A$ and $B$, such that $A$ and $B$ are derived Morita equivalent if and only if $G \cong 1$.

\subsection*{2. Simplicial Construction and Homotopy Type}
We construct a finite connected simplicial set $S$ with a single vertex $v$ (i.e., $S_0 = \{v\}$). The non-degenerate simplices of dimensions 1 and 2 are defined as follows:
\begin{itemize}
    \item \textbf{1-simplices:} For each generator $i \in \{1, \dots, n\}$, we introduce two 1-simplices, $x_i$ and $y_i$ (representing $x_i$ and $x_i^{-1}$). For each relation $r_j = z_{j,1} z_{j,2} \dots z_{j,L_j}$ (where each $z_{j,k} \in \{x_i, y_i\}$), we introduce $L_j - 2$ internal 1-simplices $e_{j,1}, \dots, e_{j,L_j-2}$. By our assumption that $L_j \ge 3$, this set of internal edges is strictly non-empty and well-defined, eliminating any boundary index contradictions. In total, $S$ has $c_1 = 2n + L - 2m$ non-degenerate 1-simplices.
    \item \textbf{2-simplices:} For each $i \in \{1,\dots,n\}$, we introduce two 2-simplices $\sigma_{i,1}$ and $\sigma_{i,2}$ to enforce inversion: $\sigma_{i,1}$ has boundaries $\partial_2 = x_i, \partial_0 = y_i, \partial_1 = s_0v$ (the degenerate 1-simplex); and $\sigma_{i,2}$ has $\partial_2 = y_i, \partial_0 = x_i, \partial_1 = s_0v$.
    For each relation $r_j$, we introduce a sequence of $L_j - 1$ 2-simplices to concatenate the word: $\tau_{j,1}$ has $\partial_2 = z_{j,1}, \partial_0 = z_{j,2}, \partial_1 = e_{j,1}$. For $2 \le k \le L_j - 2$, $\tau_{j,k}$ has $\partial_2 = e_{j,k-1}, \partial_0 = z_{j,k+1}, \partial_1 = e_{j,k}$. Finally, $\tau_{j,L_j-1}$ has $\partial_2 = e_{j,L_j-2}, \partial_0 = z_{j,L_j}, \partial_1 = s_0v$. In total, $S$ has $c_2 = 2n + L - m$ non-degenerate 2-simplices.
\end{itemize}
All simplices of dimension $\ge 3$ are set to be degenerate.

The fundamental group of the geometric realization $|S|$ corresponds to the edge monoid $M$ generated by the 1-simplices, modulo the relations $\partial_1 \sigma = (\partial_2 \sigma)(\partial_0 \sigma)$ for each 2-simplex $\sigma$. The $\sigma$-simplices enforce $x_i y_i = 1$ and $y_i x_i = 1$, strictly making every $x_i$ and $y_i$ invertible. The $\tau$-simplices define the intermediate edges as $e_{j,k} = z_{j,1} \dots z_{j,k+1}$ and ultimately enforce $r_j = 1$. Because all generators of $M$ are invertible, $M$ is explicitly a group, and by construction $M \cong G$.

Topologically, $|S|$ is a finite CW complex. By sequentially collapsing the internal edges $e_{j,k}$ along the 2-simplices $\tau_{j,k}$, each sequence of $L_j - 1$ 2-simplices coalesces into a single 2-cell attached along the loop corresponding to the word $r_j$. Using $\sigma_{i,1}$ to deformation retract $y_i$ onto $x_i^{-1}$, these $m$ amalgamated 2-cells precisely form the standard 2-dimensional presentation complex $X_P$. The remaining $n$ 2-cells $\sigma_{i,2}$ become attached along strictly null-homotopic loops $x_i^{-1} x_i$. Attaching a 2-cell along a constant loop wedges a 2-sphere. Therefore, $|S| \simeq X_P \vee \bigvee^n S^2$.

If $G \cong 1$, the complex $X_P$ is a simply connected finite 2-complex. By classical homotopy theory, it is homotopy equivalent to a wedge of 2-spheres. Computing Euler characteristics yields $\chi(X_P) = 1 - n + m$, meaning $X_P \simeq \bigvee^{m-n} S^2$. It follows that when $G \cong 1$, $|S| \simeq \bigvee^m S^2$.

\subsection*{3. Construction of the DGAs $A$ and $B$}
Let $N_*(S; R)$ be the normalized simplicial chain complex of $S$ over $R$. We define $A = \Omega N_*(S; R)$, the unreduced Adams cobar construction. Since $S$ has no non-degenerate simplices above dimension 2, $A$ is precisely the tensor algebra $T_R(N_1(S)[-1] \oplus N_2(S)[-1])$. Consequently, $A$ is the free non-commutative graded $R$-algebra on $c_1$ generators of degree 0 and $c_2$ generators of degree 1.
We define $B$ to be the free non-commutative graded $R$-algebra on $m$ homogeneous generators of degree 1, equipped with the zero differential. Both $A$ and $B$ perfectly satisfy condition (a) of the problem.

We now compute $H_0(A)$. The degree 0 part of $A$ is the free algebra $A_0 = T_R(N_1(S)[-1])$. Under standard sign conventions, the differential $d: A_1 \to A_0$ evaluates on a generator $\sigma \in N_2(S)[-1]$ as $d(\sigma) = \partial_1 \sigma - \partial_0 \sigma - \partial_2 \sigma - \partial_2 \sigma \otimes \partial_0 \sigma$. Let $F = R\langle g_e \mid e \in N_1(S) \rangle$ be the free $R$-algebra on symbols $g_e$. The change of basis $\tilde{\phi}: A_0 \to F$ given by $e \mapsto g_e - 1$ is an isomorphism of free algebras. Under $\tilde{\phi}$, we evaluate $d(\sigma)$ noting that degenerate simplices like $s_0v$ correspond to $0$ in $N_1(S)$, hence $g_{s_0v} = 1$:
$$ \tilde{\phi}(d(\sigma)) = (g_{\partial_1 \sigma} - 1) - (g_{\partial_0 \sigma} - 1) - (g_{\partial_2 \sigma} - 1) - (g_{\partial_2 \sigma} - 1)(g_{\partial_0 \sigma} - 1) = g_{\partial_1 \sigma} - g_{\partial_2 \sigma} g_{\partial_0 \sigma}. $$
Since $H_0(A) = A_0 / \langle d(A_1) \rangle$, the map $\tilde{\phi}$ descends to an isomorphism $H_0(A) \cong F / \langle g_{\partial_1 \sigma} - g_{\partial_2 \sigma} g_{\partial_0 \sigma} \rangle$. The latter quotient is identically the monoid algebra $R[M]$. Since we rigorously established $M \cong G$, we obtain $H_0(A) \cong R[G]$.

\subsection*{4. Derived Morita Equivalence requires $G = 1$}
By a fundamental theorem of B.\ Keller (\textit{Invariance and localization for cyclic homology of DG algebras}, J. Pure Appl. Algebra, 123 (1998), 223--273), if two DGAs are derived Morita equivalent, their Hochschild homologies are isomorphic. For any DGA supported in non-negative homological degrees, $HH_0$ is canonically $H_0 / [H_0, H_0]$.
For $B$, since $H_0(B) = R$, we have $HH_0(B) \cong R$.
For $A$, $HH_0(A) \cong R[G] / [R[G], R[G]] \cong \bigoplus_{[g] \in \text{Conj}(G)} R$, which is the free $R$-module generated by the conjugacy classes of $G$.

If $A$ and $B$ are derived Morita equivalent, the isomorphism $HH_0(A) \cong HH_0(B)$ forces the rank of the $R$-module of conjugacy classes to be exactly 1. Because the identity element always forms its own conjugacy class, $G$ can possess no other conjugacy classes. This strictly implies $G \cong 1$.

\subsection*{5. Sufficiency when $G = 1$}
Conversely, assume $G \cong 1$. We must prove $A$ and $B$ are derived Morita equivalent. Because the simplicial set $S$ possesses non-degenerate 1-simplices, it is not 1-reduced. Thus, we invoke the extension of Adams' cobar equivalence to connected spaces without 1-reduction assumptions (see M.\ Rivera and M.\ Zeinalian, \textit{Cubical rigidification, the cobar construction, and the based loop space}, Algebr. Geom. Topol., 18 (2018), 3789--3820). This theorem ensures that the natural DGA morphism from the unreduced cobar construction $A = \Omega N_*(S; R)$ to the singular chains on the Moore based loop space $C_*(\Omega|S|; R)$ is a quasi-isomorphism.

Because $G \cong 1$ implies $|S| \simeq \bigvee^m S^2$, we have a sequence of weak equivalences of spaces and thus DGA quasi-isomorphisms:
$$ A \;\simeq\; C_*(\Omega|S|; R) \;\simeq\; C_*\left(\Omega \Big(\bigvee^m S^2\Big); R\right). $$
By the Bott-Samelson theorem, the homology algebra $H_*\big(\Omega (\bigvee^m S^2); R\big) \cong T_R(R^m[-1])$, which as a graded algebra is exactly isomorphic to $B$. Because $B$ is a free DGA with zero differential, we can define a DGA morphism by mapping its $m$ generators to cycle representatives of a basis for $H_1\big(\Omega (\bigvee^m S^2); R\big)$. This unique lift yields a DGA quasi-isomorphism $B \xrightarrow{\simeq} C_*\big(\Omega (\bigvee^m S^2); R\big)$.

Thus, $A$ and $B$ are connected by a zig-zag of quasi-isomorphisms, establishing that they are derived Morita equivalent. Since the derived Morita equivalence of these finitely generated free DGAs correctly decides the triviality of a finitely presented group, the problem is undecidable.
\end{solution}



\begin{solution}{Another model response (\emph{Aletheia})}
\textbf{Theorem.} \textit{Let $R$ be $\mathbb{Z}$ or a field. The problem of determining whether two differential graded algebras of type (a) are derived Morita equivalent is algorithmically undecidable. Consequently, the general decision problem is undecidable.}

\begin{proof}
The proof proceeds in four steps.

\textbf{Step 1: The Triviality Problem for Finitely Presented Algebras}

By the foundational Markov-Post theorem (A. A. Markov, 1947; E. L. Post, 1947), the word problem for finitely presented monoids is algorithmically undecidable. Specifically, there exists a fixed finitely presented monoid $S = \langle x_1, \dots, x_n \mid u_1 = v_1, \dots, u_m = v_m \rangle$ such that there is no general algorithm to decide whether two arbitrary input words $w_1$ and $w_2$ represent the same element in $S$.

Given an input pair of words $(w_1, w_2)$, we construct a finitely presented associative $R$-algebra $P$. The generators of $P$ are the monoid generators $x_1, \dots, x_n$, plus two new generators $y$ and $z$. The relations defining $P$ are:
\begin{enumerate}
    \item $u_j = v_j$ for $j = 1, \dots, m$.
    \item $y(w_1 - w_2)z = 1$.
\end{enumerate}

We claim that $P$ is the zero ring ($P = 0$, meaning $1 = 0$ in $P$) if and only if $w_1 = w_2$ in $S$.

\begin{itemize}
    \item \textbf{Suppose $w_1 = w_2$ in $S$:} Then the monoid relations strictly enforce $w_1 - w_2 = 0$ identically in the algebra $P$. Substituting this into the final relation yields $y(0)z = 1 \implies 0 = 1$. Thus, $P$ algebraically collapses to the zero ring.

    \item \textbf{Suppose $w_1 \neq w_2$ in $S$:} We construct an explicit, non-trivial module representation of $P$. Let $A = R[S]$ be the monoid algebra of $S$ over $R$. Because distinct elements of $S$ form a canonical free $R$-basis for $A$, the distinct words $w_1$ and $w_2$ are linearly independent, meaning $w_1 - w_2 \neq 0$ in $A$. We define an $R$-linear functional $f: A \to R$ on the basis $S$ by setting $f(w_1) = 1$ and $f(s) = 0$ for all $s \in S \setminus \{w_1\}$. This ensures $f(w_1 - w_2) = 1 - 0 = 1$.

    Let $V = \bigoplus_{k=1}^\infty A^{\otimes_R k}$ be the strictly positive-degree components of the tensor algebra of $A$ over $R$, viewed as an $R$-module. We denote the degree $k$ components by $V_k$. Canonically, $V_1 = A$.

    We define the action of the generators of $P$ as $R$-linear endomorphisms on $V$. We equip $V$ with a left $A$-module structure by defining the action $L_a : V \to V$ for any $a \in A$ strictly for $k \ge 1$:
    \[
        L_a(r_1 \otimes \dots \otimes r_k) = (a r_1) \otimes \dots \otimes r_k \in V_k.
    \]
    Because $L_{1_A}(r_1 \otimes \dots) = (1_A r_1) \otimes \dots = r_1 \otimes \dots$, we strictly have $L_{1_A} = \operatorname{id}_V$. This makes $V$ a valid unital left $A$-module, ensuring all monoid relations $u_j = v_j$ (even those equating to the empty word) are exactly satisfied when mapped into $\operatorname{End}_R(V)$.

    We define the action of $z$ as the $R$-linear shift operator $Z: V \to V$, mapping $Z(v) = 1_A \otimes v$.

    We define the action of $y$ as the $R$-linear operator $Y: V \to V$, which acts as $0$ on the degree $1$ component $V_1 = A$. For $k \ge 2$, we identify $V_k \cong A \otimes_R V_{k-1}$ and define $Y: V_k \to V_{k-1}$ as the evaluation map $f \otimes \operatorname{id}_{V_{k-1}}$, yielding $Y(r \otimes v) = f(r)v$.

    To verify the final relation, consider the operator $Y \circ L_{w_1 - w_2} \circ Z$ applied to any $v \in V_k$ for $k \ge 1$:
    \begin{align*}
        Z(v) &= 1_A \otimes v \in V_{k+1}\\
        L_{w_1 - w_2} (1_A \otimes v) &= (w_1 - w_2) \otimes v \in V_{k+1}\\
        Y((w_1 - w_2) \otimes v) &= f(w_1 - w_2) v = 1 \cdot v = v \in V_k.
    \end{align*}
    Thus, $Y \circ L_{w_1 - w_2} \circ Z = \operatorname{id}_V$ unconditionally. Since all relations of $P$ are perfectly satisfied, there is a well-defined $R$-algebra homomorphism $P \to \operatorname{End}_R(V)$. Because the monoid $S$ contains at least the empty word, $V_1 = A \neq 0$, meaning the identity map $\operatorname{id}_V \neq 0$. Therefore, $1 \neq 0$ in $P$, meaning $P \neq 0$.
\end{itemize}

Because deciding whether $w_1 = w_2$ in $S$ is algorithmically undecidable, the triviality problem for the finitely presented associative $R$-algebra $P$ is undecidable.

\textbf{Step 2: Construction of the DGA $\mathcal{A}$}

Let the finite presentation of the algebra from Step 1 be rewritten abstractly as $P = R\langle X_1, \dots, X_N \rangle / (r_1, \dots, r_M)$. We computably construct a DGA $\mathcal{A}$ of type (a) using cohomological grading. Let $\mathcal{A}$ be the free non-commutative graded $R$-algebra on $N+M$ homogeneous generators defined as follows:
\begin{itemize}
    \item Generators $X_1, \dots, X_N$ in degree $0$.
    \item Generators $Y_1, \dots, Y_M$ in degree $-1$.
\end{itemize}

We define the differential $d$ of degree $+1$ on the generators by:
\begin{align*}
    d(X_i) &= 0 \quad \text{for all } i = 1, \dots, N, \\
    d(Y_j) &= r_j(X_1, \dots, X_N) \quad \text{for all } j = 1, \dots, M.
\end{align*}

We extend $d$ to all of $\mathcal{A}$ via the standard graded Leibniz rule $d(ab) = d(a)b + (-1)^{|a|}a d(b)$. Because the relations $r_j$ are polynomials exclusively in the degree $0$ generators $X_i$, the derivation property ensures $d(r_j) = 0$. Thus, $d^2(Y_j) = d(r_j) = 0$. By evaluating $d^2(ab)$, one easily verifies that $d^2$ acts globally as an even derivation (degree $+2$) without parity signs: $d^2(ab) = d^2(a)b + a d^2(b)$. An even derivation that vanishes on all free generators must vanish globally, so $d^2 = 0$. This makes $\mathcal{A}$ a valid DGA strictly of type (a).

\textbf{Step 3: Acyclicity of $\mathcal{A}$}

Because $\mathcal{A}$ is freely generated and has no generators in degrees less than $-1$, any element $u \in \mathcal{A}^{-1}$ is a linear combination of monomials containing exactly one generator of degree $-1$. Thus, $u$ takes the form
\[
    u = \sum_k c_k p_k Y_{j_k} q_k,
\]
where $c_k \in R$ and $p_k, q_k$ are monomials strictly in the $X_i$'s (meaning $|p_k| = 0$).

Applying the graded Leibniz rule:
\[
    d(p_k Y_{j_k} q_k) = d(p_k) Y_{j_k} q_k + (-1)^{|p_k|} p_k d(Y_{j_k}) q_k + (-1)^{|p_k| + |Y_{j_k}|} p_k Y_{j_k} d(q_k).
\]
Because $d(X_i) = 0$, we have $d(p_k) = d(q_k) = 0$. The parity sign evaluates to $(-1)^{|p_k|} = (-1)^0 = 1$, yielding $d(p_k Y_{j_k} q_k) = p_k r_{j_k} q_k$. Extending linearly gives:
\[
    d(u) = \sum_k c_k p_k r_{j_k} q_k.
\]
This demonstrates that the image of the differential $d(\mathcal{A}^{-1}) \subset \mathcal{A}^0$ precisely coincides with the two-sided ideal $I = (r_1, \dots, r_M)$ generated by the relations. Consequently, $1 \in \mathcal{A}^0$ is an exact boundary ($1 = d(u)$ for some $u \in \mathcal{A}^{-1}$) if and only if $1 \in I$, which holds if and only if $P = 0$.

Suppose $P = 0$. Then $1 = d(u)$ for some $u \in \mathcal{A}^{-1}$. The element $u$ is intrinsically homogeneous of degree $-1$ ($|u| = -1$). For \textit{any} arbitrary cycle $c \in \mathcal{A}$ (meaning $d(c) = 0$), left-multiplication by $u$ operates as a global contracting homotopy. Because $u$ is homogeneous, the graded Leibniz rule extends cleanly regardless of the homogeneity of $c$:
\[
    d(u c) = d(u)c + (-1)^{|u|}u d(c) = 1 \cdot c + (-1)^{-1}(u \cdot 0) = c.
\]
This proves every cycle is a boundary, meaning $\mathcal{A}$ is exact and therefore acyclic ($\mathcal{A} \simeq 0$).
Conversely, if $\mathcal{A} \simeq 0$, the zeroth cohomology $H^0(\mathcal{A}) = 0$. The multiplicative identity $1 \in \mathcal{A}^0$ is trivially a cycle ($d(1) = 0$), so it must be a boundary, implying $1 \in I$ and hence $P = 0$. Thus, $\mathcal{A} \simeq 0 \iff P = 0$.

\textbf{Step 4: The Target DGA $\mathcal{B}$ and Conclusion}

Let $\mathcal{B}$ be the fixed DGA of type (a) defined as the free non-commutative $R$-algebra on a single generator $Y$ in degree $-1$, equipped with the differential $d(Y) = 1$. Because $1$ is an exact boundary in $\mathcal{B}$, applying the exact same contracting homotopy logic $c = d(Yc)$ ensures that $\mathcal{B}$ is unconditionally acyclic ($\mathcal{B} \simeq 0$).

By foundational derived Morita theory for DG algebras (B. Keller, ``Deriving DG categories'', \textit{Ann. Sci. École Norm. Sup.}, 1994), two DGAs $\mathcal{X}$ and $\mathcal{Y}$ are derived Morita equivalent if and only if their derived categories of DG modules $\mathcal{D}(\mathcal{X})$ and $\mathcal{D}(\mathcal{Y})$ are equivalent as triangulated categories. For any DGA $\mathcal{X}$, the free module $\mathcal{X}$ is a compact generator for $\mathcal{D}(\mathcal{X})$, with natural abelian group isomorphisms $\operatorname{Hom}_{\mathcal{D}(\mathcal{X})}(\mathcal{X}, \mathcal{X}[n]) \cong H^n(\mathcal{X})$. Therefore, $\mathcal{D}(\mathcal{X})$ is equivalent to the zero category $0$ if and only if its compact generator is acyclic ($H^\bullet(\mathcal{X})=0$).

Since $\mathcal{B} \simeq 0$, its derived category $\mathcal{D}(\mathcal{B}) \simeq 0$. Thus, $\mathcal{A}$ is derived Morita equivalent to $\mathcal{B}$ if and only if $\mathcal{D}(\mathcal{A}) \simeq 0$, which holds if and only if $\mathcal{A} \simeq 0$.

Tracing back the logic mappings:
\[
    \mathcal{A} \text{ is derived Morita equivalent to } \mathcal{B} \iff \mathcal{A} \simeq 0 \iff P = 0 \iff w_1 = w_2.
\]

Because the target DGA $\mathcal{B}$ is fixed, and the DGA $\mathcal{A}$ is deterministically generated from the inputs $w_1, w_2$ of the monoid $S$, any general algorithm deciding derived Morita equivalence for DGAs of type (a) would successfully decide the word problem for $S$. Since the latter is undecidable, there can be no general algorithm to determine whether $\mathcal{A}$ is derived Morita equivalent to $\mathcal{B}$. The overarching problem is unconditionally algorithmically undecidable.
\end{proof}
\end{solution}


\end{document}
